
\section{个人收获与工程能力体现}

通过本课题，完整经历了需求分析、架构设计、模块实现、测试与文档化流程。OpenSpec 的使用使本人认识到\textbf{可验证需求}对控制范围蔓延的重要性；ECS 实践加深了对“组合式实体”的理解；H5 性能优化让人意识到主线程预算的硬约束\cite{kw_project_management_agile,kw_ecs_game_architecture}。

在团队协作场景下（若有多人参与），\texttt{WORKSTREAMS.md} 划分的元游戏/战斗/UI 工作流可避免合并冲突；单人场景下亦可作为任务看板。版本控制 Git 与变更提案目录结合，形成“代码 + 规格”双轨历史。

\section{与模板要求的三项补充分析}

\subsection{社会、健康、安全、法律与文化影响}

游戏作为文化产品，承担娱乐与情绪调节功能，亦可能带来沉迷风险\cite{kw_ethics_game_design}。本项目为教学演示，未内置付费与社交诱导机制，降低经济纠纷风险。战斗表现采用抽象怪物与血条，避免血腥写实。建议在用户手册中提示合理作息。若面向未成年人发布，须遵守防沉迷与时长限制政策。

\subsection{绿色节能与兼容性}

低 CPU 占用意味着更低能耗与散热\cite{kw_green_web_perf}。对象池、逻辑暂停、Worker 卸载均有助于降低平均功耗。兼容性上，Worker 回退保证旧浏览器可运行；配表双通道加载适配不同发布路径；2D 渲染对 GPU 要求低于 3D 大作。未来可增加“节能模式”限制帧率与同屏怪数。

\subsection{项目管理应用心得}

采用“提案评审 $\rightarrow$ 任务清单 $\rightarrow$ 实现 $\rightarrow$ 验证 $\rightarrow$ 归档”流程，等价于短周期迭代\cite{kw_project_management_agile}。每个 \texttt{change-id} 范围可控，适合毕业设计 3--4 个月周期。风险是提案过多导致集成延迟——缓解方式是按 Must/Should 优先级排期（见第~\ref{chap:requirement}~章 MoSCoW 表）。

\section{对后续学习者的建议}

\begin{enumerate}[label=\arabic*.]
  \item 先实现 ECS 核心与单怪物追逐，再叠加子弹与配表，避免一次堆满系统；
  \item 尽早建立配表同步脚本，减少“能编不能跑”时间；
  \item 性能优化基于数据：先测量再池化/Worker，避免过早优化；
  \item 论文与 OpenSpec 同步更新，答辩前运行 \texttt{openspec validate}；
  \item 参考文献勿编造，按 \texttt{references.bib} 关键词检索权威来源。
\end{enumerate}

\section{结束语}

Survivor And Fight 作为本科毕业设计，在有限周期内探索了现代 H5 游戏客户端的一种可行架构路径。论文初稿以 LaTeX 呈现，文献以关键词占位，待作者补全后即为合格终稿基础。期待后续在法杖编程、自动化测试与性能数据方面继续完善，使项目从“可演示原型”成长为“可发布小游戏”。
